\icode{N22J-O-CA-L3-261003-003}
\nceeID[8]{2022河北}{16(1)}
如图，一个半径为\(R\)的玻璃球，\(O\)点为球心。球面内侧单色点光源\(S\)发出的一束光从\(A\)点射出，出射光线\(AB\)与球直径\(SC\)平行，\(\theta =\) \,\qty{30}{\degree}。光在真空中的传播速度为\(c\)。求：
\begin{enumerate}
	\item 玻璃的折射率；
	\item 从\(S\)发出的光线经多次全反射回到\(S\)点的最短时间。
\end{enumerate}

\begin{tikzpicture}[scale=.8,decoration={markings,mark=at position 0.5 with {\arrow{>}}},font=\scriptsize]
	\draw (2,0) circle(2);
	\draw[postaction={decorate},purple] (0,0)--(30:{2*sqrt(3)});
	\draw[postaction={decorate},purple] (30:{2*sqrt(3)})node[above,black]{\(A\)}--(5,{sqrt(3)})node[right,black]{\(B\)};
	\draw[dashed] (0,0)node[left]{\(S\)}--(4,0)node[right]{\(C\)};
	\draw (.4,0) arc(0:30:.4);
	\fill (2,0)node[below]{\(O\)} circle(.06);
	\node at(.6,.15){\(\theta\)};
\end{tikzpicture}
\begin{answer}
    （1）\(\sqrt{3}\)（2）\(\frac{4 \sqrt{6}R}{c}\)\qquad 提示：发生全反射回到\(S\)点正四边形及以上的图形，最短时间为正四边形的时间

    \begin{tikzpicture}[scale=.8,decoration={markings,mark=at position 0.5 with {\arrow{>}}},font=\scriptsize]
	\draw (2,0) circle(2);
	\draw[postaction={decorate},purple] (0,0)--(30:{2*sqrt(3)});
	\draw[postaction={decorate},purple] (30:{2*sqrt(3)})node[above,black]{\(A\)}--(5,{sqrt(3)})node[right,black]{\(B\)};
	\draw[dashed] (0,0)node[left]{\(S\)}--(4,0)node[right]{\(C\)};
	\draw (.4,0) arc(0:30:.4);
    \draw[dashed,cyan] (2,0)--++(60:3);
	\fill (2,0)node[below]{\(O\)} circle(.06);
	\node at(.6,.15){\(\theta\)};
    \end{tikzpicture}
    \begin{tikzpicture}[scale=.8,decoration={markings,mark=at position 0.5 with {\arrow{>}}},font=\scriptsize]
	\draw (2,0) circle(2);
	\draw[postaction={decorate},purple] (0,0)--(2,2);
    \draw[postaction={decorate},purple] (2,2)--(4,0);
    \draw[postaction={decorate},purple] (4,0)--(2,-2);
    \draw[postaction={decorate},purple] (2,-2)--(0,0);
	\draw[dashed] (0,0)node[left]{\(S\)}--(4,0)node[right]{\(C\)};
	\fill (2,0)node[below]{\(O\)} circle(.06);
    \end{tikzpicture}
\end{answer}